% Original numbered item 2
\item Put $a=|\lambda_*|$.  The normalized root coordinates are
\[
 z_\infty=(a/\lambda)^{1/5},\qquad
 z_0=(\lambda/a)^{1/4}.
\]
In the regular root frames of P1.16 their connections have the form
\[
 G(z)=\frac{\sum_{j=0}^{m-1}B_jz^j}{1+z^m},\qquad m=5,4.
\tag{1.18.2}
\]
The entries of $B_j$ are the exact algebraic numbers obtained by
substituting the positive real number $a^{\pm1/m}$ in the rational
connection.  The generator of the finite deck group acts on these
frames by their stated eigencharacters.  Since its flat matrix is
real and fixes $e$, it leaves $h_b$ invariant.  It therefore suffices
to cover the sector $0\leq\arg z\leq2\pi/m$ in each root disc.


% Original numbered item 3
\item Use the positive diagonal weights
\[
 (1,1/128,1/128,1/32,1/64,1/32)
\]
and the corresponding weighted maximum row norm.  Exact ball
arithmetic gives
\[
 \left\|\sum_j |B_j|\right\|<1.447\quad(m=5),\qquad
 \left\|\sum_j |B_j|\right\|<1.135\quad(m=4).
\]
Here the absolute values are entrywise, before taking the row norm.
Let $B$ denote the actual enclosed bound.  The scalar inequality
\[
 \frac1{1-t^m}\leq\frac1{m(1-t)}+\frac{m-1}{m},\qquad0\leq t<1,
\]
follows on multiplying by $m(1-t)(1+\cdots+t^{m-1})$;
the remaining assertion is
$\sum_{j=1}^{m-1}t^j\geq(m-1)t^m$.
Gronwall's inequality consequently bounds the fundamental solution
$U(0)=I$ on $|z|\leq R<1$ by
\[
 K_R=\exp\!\left[B\left(-\frac{\log(1-R)}m
                    +\frac{m-1}{m}R\right)\right].
\]
Cauchy's estimate bounds the tail after degree $n$ on $|z|\leq r<R$
by
\[
 K_R\frac{(r/R)^{n+1}}{1-r/R}.
\tag{1.18.3}
\]
The code uses $R=0.999$.  The coefficients are computed directly
from (1.18.2):
\[
 nU_n=-(n-m)U_{n-m}
       +\sum_{j=0}^{\min(m-1,n-1)}U_{n-1-j}B_j,
\]
with the first term omitted when $n<m$.
There is no derivative approximation in this recursion.
The regular limit frame is determined by certified continuation
from $\lambda=1$ to a positive root-chart point and inversion of
this convergent series.  At the infinity center it yields
\[
 h_b(\infty)\in
 [-1.094540250311402780116268263054945704
       \ \mathbin{\pm}\ 3\cdot10^{-37}].
\tag{1.18.4}
\]
The corresponding enclosure at zero has center
\[
-0.9317933763565235054307951057960111149908282005165194125
\]
and radius $7.91\cdot10^{-56}$.


% Original numbered item 5
\item For a polynomial $p(z)=\sum c_jz^j$ with ball coefficients, evaluation
on a rectangle centered at $z_c$, of enclosing radius $\rho$, uses
its midpoint-coefficient polynomial at the point $z_c$.  Coefficient
uncertainty is bounded by
$\sum_j\operatorname{rad}(c_j)|z_c|^j$, and variation over the
rectangle by
$\rho\sum_j j|c_j|r^{j-1}$, where $r$ bounds the modulus on it.
The tail (1.18.3) is then added.  The point evaluation uses 800-digit
ball arithmetic; its rounding radius is retained.  These explicit
bounds avoid the excessive interval enlargement of long rectangular
Horner recursions.  No error radius is discarded.


% Original numbered item 6
\item On the square $|\operatorname{Re}z_\infty|,
|\operatorname{Im}z_\infty|\leq0.1$, differentiating the exact linear
system for the real intersection vector gives the following
certified bounds for the real Hessian of $h_b$:
\[
\begin{split}
 0.2898644&<h_{xx},h_{yy}<0.3946911,\\
 |h_{xy}|&<0.0524256,\qquad
 \det\operatorname{Hess}(h_b)>0.0757786.
\end{split}
\tag{1.18.5}
\]
The derivatives of the period frame are evaluated by
$Y'=YG$ and $Y''=Y(G^2+G')$; the derivatives of the intersection
vector follow by differentiating its invertible four-by-four real
system.  Finite rotation invariance forces $dh_b(0)=0$ in this
root chart.  Positive definiteness on the convex square therefore
proves a strict minimum at its center.


% Original numbered item 8
\item It remains to cover a cusp neighborhood and the collar between these
two root discs.  Write
\[
 q=(\lambda-\lambda_*)/\lambda_*.
\]
In this coordinate the connection is $R/q+A(q)$, with $R^2=0$
and $A$ analytic on $|q|<1$.  A regular-singular factorization is
\[
 Y(q)=H\exp((\log q)R)P(q),\qquad P(0)=I.
\]
The recursion for its analytic factor is
\[
 (n-\operatorname{ad}_R)P_n
   =\sum_{j=0}^{n-1}P_{n-1-j}A_j,
 \qquad \operatorname{ad}_R(Q)=QR-RQ.
\]
Because $R^2=0$, the cube of this adjoint operator vanishes and its
inverse is the finite expression
$1/n+\operatorname{ad}_R/n^2+\operatorname{ad}_R^2/n^3$.
In an exact dyadic diagonal norm the code proves
$\|R\|<0.841$ and uses the sharper recorded bound
\[
 C=1+2\|R\|+4\|R\|^2<5.508.
\]
On the circle of radius $0.4$, the absolute-coefficient majorant of
$A$ is less than $4.895$.  The coefficient recursion is dominated
by the scalar equation $p'=Ca_{\rm abs}p$; hence
$\|P\|\leq\exp(C\cdot4.895\cdot0.4)$ there.  Cauchy's estimate
now gives the same explicit geometric tail bound as (1.18.3), with
this exponential in place of $K_R$ and $R=0.4$.


% Original numbered item 10
\item The only remaining region, expressed in the infinity-root chart, is
\[
 0.98\leq |z_\infty|\leq0.98^{-4/5},\qquad
 0\leq\arg z_\infty\leq2\pi/5,
 \qquad |q|\geq0.25,
\]
where $q=-z_\infty^{-5}-1$.
Each rectangle in this collar is reached from radius $0.95$ by
Taylor continuation of (1.18.2).  The local analytic radius is
$0.449$ times the distance to the nearest root of $1+z^5$.
The numerator is bounded in the same diagonal norm and the
denominator by the product of the five distances to those roots.
Gronwall and Cauchy give an explicit tail for every step, exactly
as in P1.16.  There are 427 steps and 117 accepted rectangles,
of maximum depth eight; each proves $h_b>h_b(\infty)$.

